% from ega4-18-fr.tex
\oldpage[IV-4]{199}

et si $Y$ est régulier, alors il existe un voisinage ouvert $U$ de $Y$ dans $X$ tel que $X$ soit régulier aux points de $U-Y$.
\end{corollary}

Pour prouver la première assertion, on peut se borner au cas où $X$ est un schéma local de point fermé $x$; l'hypothèse entraîne alors que $Y$ est régulier ($0$, 17.3.2), donc il en est de même de $Y'$ par (19.4.4). Notons maintenant que le morphisme $f:X'\to X$ est projectif, donc propre, et par suite (II, 7.2.1), si l'on considère l'unique point $x'_1$ de $X'$ au-dessus de $x_1$ (II, 8.1.3), il existe une spécialisation de $x'_1$ appartenant à $Y'$; en lui appliquant (19.4.4) et utilisant ($0$, 17.3.2) on en conclut que $X'$ est régulier au point $x'_1$, donc $X$ régulier au point $x_1$ puisque $X'-Y'$ est isomorphe à $X-Y$ (II, 8.1.3). Si $Y$ est régulier, tout point de $X-Y$ qui est générisation d'un point de $Y$ appartient donc à $\operatorname{Reg}(X)$; si $\operatorname{Reg}(X)$ est ouvert, $U=Y\cup\operatorname{Reg}(X)$ est donc localement constructible et stable par générisation, donc ouvert ($0_{\mathrm{III}}$, 9.2.5) et répond à la question.

\begin{proposition}[19.4.6]
\label{IV.19.4.6}
Soient $g:X\to S$ un morphisme, $Y$ un sous-préschéma fermé de $X$ défini par un Idéal quasi-cohérent $\mathscr J$, $X'$ le $X$-schéma obtenu en faisant éclater $Y$, $Y'$ l'image réciproque de $Y$ dans $X'$.
\begin{enumerate}
  \item[(i)] Les conditions suivantes sont équivalentes:
  \begin{enumerate}
    \item[a)] $\operatorname{gr}_{\mathscr J}^*(\sh{O}_X)$ est une $\sh{O}_X$-Algèbre $g$-plate.
    \item[b)] Pour tout $n\geq0$, le $\sh{O}_X$-Module $\sh{O}_X/\mathscr J^{n+1}$ est $g$-plat (en d'autres termes, le $n$-ème voisinage infinitésimal $Y^{(n)}$ de $Y$ dans $X$ (16.1.2) est $S$-plat).
    \item[c)] Pour tout changement de base $S_1\to S$ et tout entier $n\geq1$, si l'on pose $X_1=X\times_S S_1$, l'homomorphisme canonique $\mathscr J^n\otimes_{\sh{O}_S}\sh{O}_{S_1}\to\mathscr J^n\sh{O}_{X_1}$ est bijectif.
  \end{enumerate}
  Lorsqu'il en est ainsi, $Y'$ est $S$-plat, et pour tout changement de base $S_1\to S$, si $Y_1$ est l'image réciproque de $Y$ dans $X_1$, le préschéma $X'_1$ obtenu en faisant éclater $Y_1$ dans $X_1$ est canoniquement isomorphe à $X'\times_S S_1$.
  \item[(ii)] Supposons que les conditions équivalentes de (i) soient satisfaites et de plus que les morphismes $X\to S$ et $Y\to S$ soient localement de présentation finie. Alors $\mathscr S=\bigoplus_{n\geq0}\mathscr J^n$ est une $\sh{O}_X$-Algèbre de présentation finie, le morphisme $X'\to X$ est de présentation finie, l'immersion canonique $Y'\to X'$ est transversalement régulière de codimension $1$ relativement à $S$, et l'ensemble des points de $X$ (resp. $X'$) où $X$ (resp. $X'$) est $S$-plat, est un voisinage de $Y$ (resp. $Y'$).
\end{enumerate}
\end{proposition}

\begin{enumerate}
  \item[(i)] L'équivalence de $a)$ et $b)$ résulte aussitôt de ($0_{\mathrm I}$, 6.1.2); pour montrer l'équivalence de $b)$ et $c)$, on peut se borner au cas où $S=\Spec(A)$ et $X=\Spec(B)$ sont affines, avec $\mathscr J=\widetilde{\mathfrak J}$, où $\mathfrak J$ est un idéal de $B$; la suite exacte des Tor montre alors que la condition $c)$ implique que $\operatorname{Tor}_1^A(B/\mathfrak J^{n+1},A')=0$ pour tout entier $n\geq0$ et toute $A$-algèbre $A'$; en prenant pour $A'$ une algèbre somme directe $A\oplus M$, où $M$ est un $A$-module quelconque et où la multiplication est définie par $(a,m)(a',m')=(aa',am'+a'm)$ on voit que la condition précédente équivaut à $\operatorname{Tor}_1^A(B/\mathfrak J^{n+1},M)=0$, donc au fait que $B/\mathfrak J^{n+1}$ est un $A$-module plat. Si l'on pose $\mathscr J_1=\mathscr J\otimes_{\sh{O}_S}\sh{O}_{S_1}$, on a alors $\bigoplus_{n\geq0}\mathscr J_1^n=(\bigoplus_{n\geq0}\mathscr J^n)\otimes_{\sh{O}_S}\sh{O}_{S_1}$ et l'assertion relative au préschéma éclaté $X'_1$ s'en déduit aussitôt. Enfin, pour prouver que $Y'$ est $S$-plat, on peut encore se borner au cas où $S$ et $X$ sont affines; si l'on pose $C=\operatorname{gr}_{\mathfrak J}^*(B)$, tout point de $Y'$ admet alors, en vertu de (19.4.2) et de (II, 2.3.6), un

